% !TEX root = ../../main.tex
%--------------------------------------------------------------
%  Chapter 5 – Limitations, Future Directions, and Conclusion
%--------------------------------------------------------------
\chapter{Limitations, Future Directions, and Conclusion}
\label{chap:conclusion}

%--------------------------------------------------------------
\section{Limitations}
\label{sec:limitations}
%--------------------------------------------------------------

Every model has constraints, and they should be stated plainly.

\textbf{Scenario Calibration vs.\ Empirical Estimation.}
The MAEI is a structured counterfactual scenario, not an econometric estimation
trained on observed labor market displacement.  Such data simply does not exist
yet; federal employment statistics have not had time to capture the effects of
technologies deployed only since 2023.  \citet{Anthropic_2026} found limited
employment effects to date, suggesting the full labor market impact of GenAI may
take years to materialize.  Until longitudinal displacement data becomes
available, the MAEI remains a ``what-if'' structural model rather than a
predictive one.

\textbf{Circular Uplift Calibration.}
The Broad Exposure Uplift weights ($W_{\text{exposure}} = 10$,
$W_{\text{protection}} = 3$) were calibrated to approximately match
\citet{Eloundou_Manning_2023}'s finding that 80\% of the workforce faces some
LLM exposure.  The sensitivity analysis (Table~\ref{tab:sensitivity}) shows that
all tested configurations produce qualitatively similar results
(82.9\%--87.5\% of occupations affected), but the specific weight values are
informed by the same literature that motivates the research question.  This
circularity does not invalidate the rankings (which are stable across weight
changes) but limits the ability to treat the absolute exposure magnitudes as
independently derived.

\textbf{Tree Saturation.}
As documented in Section~\ref{sec:tree_saturation}, the Stacked Ensemble's
tree-based components produce near-zero prediction deltas when features are
multiplied beyond the training range.  The ML model validates feature selection
and structural plausibility but does not drive the final occupation rankings.
Future work could explore whether linear or kernel-based models that can
extrapolate would produce more informative model-driven deltas, though the
meaningfulness of extrapolating far beyond training distributions is itself
debatable.

\textbf{Shared O*NET Substrate.}
Both the MAEI and \citet{Felten_Raj_Seamans_2021}'s AIOE dataset are constructed
from O*NET occupational data.  While the measurement approaches differ
substantially (scenario-based multipliers vs.\ crowdsourced AI-task linkages),
the shared underlying data source means the $\rho = +0.764$ correlation may
partially reflect the common O*NET structure rather than fully independent
methodological convergence.

\textbf{Ontological Ambiguity of ``Exposure.''}
The MAEI measures capability overlap, not job loss.  As
\citet{Acemoglu_Restrepo_2018} emphasize, the path from task automation to job
displacement runs through economic adoption costs, regulatory barriers,
organizational inertia, and worker adaptability.  A high MAEI score for
``Lawyer'' means AI can perform many legal tasks; it does not predict how many
lawyers will be employed five years from now.

\textbf{O*NET Data Latency.}
The O*NET database is updated on a rolling basis, and definitions for recently
emerged occupations (for example, ``Prompt Engineer'' or ``AI Safety
Researcher'') may not yet be captured in the current taxonomy.

%--------------------------------------------------------------
\section{Future Directions}
\label{sec:future_directions}
%--------------------------------------------------------------

Three research extensions emerge from the MAEI framework:

\textbf{Wage Shock Simulations.}
By joining the MAEI with Bureau of Labor Statistics Occupational Employment and
Wage Statistics (OEWS), researchers can calculate the total wage volume
``exposed'' to AI capability overlap.  Rather than predicting absolute job loss,
future simulations could model wage-compression scenarios: if GenAI increases
productivity in high-MAEI roles by 30\%, how does the resulting supply shock of
cognitive output depress median wages in those sectors?

\textbf{The ``Human Premium.''}
Isolating occupations with high protective feature scores (Originality, Therapy
and Counseling, Social Perceptiveness) and regressing these against longitudinal
wage data could empirically test whether demand for distinctly human skills is
increasing.  If AI commoditizes cognitive output while human authenticity,
empathy, and creative originality become scarcer, classic economic theory predicts
that wages for these human-centric skills should rise, a measurable ``Human
Premium.''

\textbf{Labor Reallocation Networks.}
Using cosine similarity between the O*NET feature vectors of high-MAEI and
low-MAEI occupations, future work could construct reskilling pathway maps.  These
algorithmic ``lifeboats'' would identify the most efficient transitions for
workers in highly exposed cognitive roles to move into roles with strong human
structural advantages, minimizing the retraining distance while maximizing the
reduction in AI exposure.

%--------------------------------------------------------------
\section{Conclusion}
\label{sec:conclusion}
%--------------------------------------------------------------

Revisiting the three research questions posed in
Chapter~\ref{chap:introduction}, the evidence supports affirmative answers to
each.  For RQ1 (structural plausibility), the Stacked Ensemble's cross-validated
$R^2$ of 0.6794 and test $R^2$ of 0.6433 demonstrate that 2013 automation
probabilities can be reliably predicted from modern O*NET features, and the
feature importance hierarchy (Fluency of Ideas at 0.2177, Originality at 0.1919)
confirms that the model learned \citet{Frey_Osborne_2017}'s theoretical
bottleneck structure rather than statistical noise.  For RQ2 (convergence with
external benchmarks), the Spearman correlation of $\rho = +0.764$
($p < 10^{-150}$) between the pure exposure delta and
\citet{Felten_Raj_Seamans_2021}'s independently constructed AIOE provides strong
evidence of construct validity.  The dual-correlation signal ($\rho = -0.357$
anchored, $\rho = +0.764$ pure delta) is a quantitative marker of the shift in risk profile.  For RQ3 (robustness to assumptions), the seven-scenario
ablation study ($\rho$ range: 0.7636--0.7651) shows that
occupation rankings are invariant to multiplier specification, converting a
potential methodological vulnerability into a demonstrated structural strength.

The Modern AI Exposure Index bridges the analytical gap between the
robotics-focused automation anxiety of the 2010s and the cognitive-AI reality of
the 2020s.  Through a five-stage computational pipeline (data engineering, NLP
feature extraction, ensemble modeling, scenario-based perturbation, and external
validation) I constructed a transparent, reproducible map of occupational
exposure to Generative AI across 1,016 U.S.\ occupations.

The external validation produced a dual-correlation signal: $\rho = +0.764$
($p < 10^{-150}$) for the pure exposure delta against
\citet{Felten_Raj_Seamans_2021}'s AIOE benchmark, and $\rho = -0.357$ for the
historically anchored score.  Together, these correlations provide evidence
of both construct validity and the structural inversion in automation risk: the
same occupations rated as safe in 2013 are now rated as highly exposed to AI
capability overlap.

The seven-scenario multiplier ablation study ($\rho$ range: 0.7636--0.7651)
proves that the occupation rankings are structurally stable, driven by the
classification of features into risk and protective categories rather than by
specific multiplier values.  The explicit decomposition of the pipeline into a
three-layer architecture (where the ML model validates feature selection, the
feature classification encodes domain knowledge, and the uplift formula generates
rankings) makes the MAEI more interpretable and defensible than a monolithic
machine learning approach.

The MAEI does not claim to predict unemployment or forecast the economic
trajectory of any specific occupation.  What it provides is a rigorous,
parameter-bounded framework for understanding which occupations overlap most with
current AI capabilities, how that profile compares to the historical baseline,
and how stable those conclusions are across a wide range of assumptions.  As
Generative AI continues to evolve, the MAEI's fully open-source pipeline can be
recalibrated with updated multipliers, extended to new O*NET releases, and
validated against emerging labor market data as it becomes available.

In the broadest terms, this dissertation contributes to a growing body of
evidence that the relationship between human labor and artificial intelligence is
entering a fundamentally new phase.  The occupations that built the middle class
of the 20th century (manufacturing, construction, transportation) were the
primary targets of the Second Wave of automation.  The occupations that promised
upward mobility in the 21st century (law, finance, programming, analysis) are now
the primary targets of the Third Wave.  The MAEI does not predict how this
transition will unfold, but it provides a structural map that policymakers,
educators, and workers can use to navigate it.  Understanding which capabilities
overlap most with AI is the essential first step toward building educational
curricula, workforce development programs, and economic policies that prepare the
labor force for a world where cognitive work is no longer exclusively human.
